\pagebreak

\chapter{Coding Recommendations}


    \section{minimising variable assignments}

        Assignments to a static or queue variable should be kept to a minimum. It
        is preferable to use different variables where possible rather than
        re-use the same variable. The more assignments a variable has the more
        delay is incurred in its input logic and hence the more likely it is
        that the logic will not meet timing constraints. In general flip-flops
        are in excess in FPGAs and it is better to use many variables with few
        assignments each rather than a few variables with many assignments each.
    
    \section{reset() better than assignment}

        It is better to use the {\bf reset()} inbuilt procedure rather than
        re-assigning the initial value (usually zero) to a static variable. The
        reset does not require the input multiplexing logic used by
        assignments, has a separate input to the flip-flops, has precedence
        over assignment in a clash should that prove useful and multiple
        resets in the code are simply ORed which is much simpler than
        assignment multiplexing. Note: the compiler now recognises when
        an assignment of a constant to a static variable returns it to its
        initial state and generates a reset instead of an assignment. In view
        of the priority of a reset over an assignment should they clash it may
        still enhance readability to use reset() explicitly.
    
    \section{queue subsets of structs}
    
        Where a processing pipeline has more fields in the queue structs as one
        progresses along the pipeline it is better to use one struct type
        containing all the necessary fields for every intermediate queue, rather
        than giving each queue its own struct, each a superset of the previous.
        This simplifies the source code while not affecting the logic resources
        used as queue (and static) elements which are not referenced are removed
        during netlist optimisation.
    
    \section{queue writes exclusivity}

        Writes to a queue must be mutually exclusive (simultaneous writes to
        component array members or fields is OK). In order to control writes
        to a queue it is usually best to restrict writing to any one queue to a
        single module.
    
    \section{nesting of target control structures}
    
        Nesting of target control structures ({\bf when}, {\bf branch}, {\bf seq
        while }, {\bf par while}, {\bf sync while}, {\bf do seq while}, {\bf do
        par while}, {\bf do sync while}) should be minimised. For example it is
        better to use independent wider {\bf when}s than to nest them. For
        example if the following is at the top level (not nested in any other
        control structures) -
        \begin{lstlisting}[gobble=8, language=threepl]
            when (l1) par {
                when (l2)
                    a <- b;
                c <- d;
            }
        \end{lstlisting}
        it would be better written as -
        \begin{lstlisting}[gobble=8, language=threepl]
           when (l1 & l2)
               a <- b;
           when (l1)
               c <- d;
        \end{lstlisting}
        Here the two {\bf when} statements are executed independently in
        parallel since they are at the top level. If the example was not
        at the top level then instead we could code it as -
        \begin{lstlisting}[gobble=8, language=threepl]
            par {
                when (l1 & l2)
                    a <- b;
                when (l1)
                    c <- d;
            }
        \end{lstlisting}

        In these sorts of transformations care must be taken to ensure that the
        two forms are equivalent in execution since statements in a {\bf par} do
        not necessarily execute the same as statements in separate top-level
        infinite loops. for example -
        \begin{lstlisting}[gobble=8, language=threepl]
            par {
                seq {
                    a <- x;
                    b <- y;
                }
                c <- z;
            }
        \end{lstlisting}
        executes differently to top-level -
        \begin{lstlisting}[gobble=8, language=threepl]
            seq {
                a <- x;
                b <- y;
            }
            c <- z;
        \end{lstlisting}
        since the {\bf seq} will take 2 clock cycles to execute and thus
        assignment to {\bf c} will occur every 2 clock cycles in the first
        example but every clock cycle in the second.
        
        As another example -
        \begin{lstlisting}[gobble=8, language=threepl]
            sync {
                when (l1)
                    a.f1 <<- 2;
                else when (l2 && l3)
                    a.f1 <<- 3;
                else when (l2 && !l3)
                    a.f1 <<- 4;
                else
                    a.f1 <<- 0;
                a.f2 <<- b;
            }
        \end{lstlisting}
        Note that a {\bf sync} is needed here to ensure that the queue is
        available to write before testing {\bf l1}, {\bf l2} etc as these
        variables could be changed by other code running in parallel while
        execution here is waiting for queue availability. This code is valid but
        will have long cascaded control logic. We could code this instead as -
        \begin{lstlisting}[gobble=8, language=threepl]
            sync {           
                when (l1 && !l2))
                    a.f1 <<- 2;
                when (!l1 && l2 && l3)
                    a.f1 <<- 3;
                when (!l1 && l2 && !l3)
                    a.f1 <<- 4;
                a.f2 <<- b;
            }
        \end{lstlisting}
        as here the {\bf when} statements execute in parallel but with a
        shallower depth of expressions to test. Note that we have been able to
        eliminate the zero assignment case because a queue assignment will occur
        because of {\bf \verb'a.f2 <<- b;'} and hence if none of the {\bf when}s
        are true, field {\bf f1} will be assigned zero by default, so the {\bf
        else} is not needed.
        
        \index{constraint!timing}
        \index{timing constraint}
        The following example illustrates ways of recoding to overcome
        timing constraint problems. An input queue carries a struct with
        two numeric fields {\bf f1} and {\bf f2} and a logical field
        {\bf coord}. If {\bf coord} is true the input queue is to be copied
        to the output queue. If {\bf coord} is false {\bf f1} is to be written
        to the output queue a number of times where the count is given by
        {\bf f2}. This can be neatly coded as -
        \begin{lstlisting}[gobble=8, language=threepl]
           struct(s,
               f1,     "uint:12",
               f2,     "uint:19",
               coord,  "log"
           );
           queue({in, out}, s);
           static(e, "uint:12");
           static(count, "uint:19", 1);
           when (in.coord)
               out <<- in;
           else seq {
               sync {
                   count <- in.f2;
                   e <- in.f1;
                   out.f1 <<- in.f1;
               }
               par while (count != 1) {
                   out.f1 <<- e;
                   count--;
               }
           }
        \end{lstlisting}
        This version is succinct but it did not meet a 10ns timing
        constraint by 5ns. The main problem is a path from static variable
        {\bf count} through the {\bf while()} test back through the containing
        when and then through to one of the assignments to {\bf out}. We can
        eliminate the {\bf while()} by realising that if we move to the top
        level of the module we have an implied infinite loop. A {\bf when()}
        at the top level therefore behaves like a {\bf while()} lower down.
        With this in mind we can recode the example as follows -
        \begin{lstlisting}[gobble=8, language=threepl]
           sync {
               when (count == 1) {
                   when (in.coord)
                       out <<- in;
                   else par {
                       count <- in.f2;
                       e <- in.f1;
                       out.f1 <<- in.f1;
                   }
               }
           }
           when (count != 1) par {
               out.f1 <<- e;
               count--;
           }
        \end{lstlisting}
        This code is perhaps less clear, but it is simpler than the previous
        version and familiarity with the implied infinite loops of 3PL reduces
        its apparent obscurity. This code still failed to meet a 10ns timing
        constraint by 1ns. The 19 bit comparison introduces a delay of 3ns or
        so and we could eliminate this by doing the comparison in parallel in
        the previous iteration and storing the result in a static variable
        as follows -
        \begin{lstlisting}[gobble=8, language=threepl]
           static(counteq1, "log", true);
           sync {
               when (counteq1) {
                   when (in.coord)
                       out <<- in;
                   else par {
                       count <- in.f2;
                       counteq1 <- (in.f2 == 1);
                       e <- in.f1;
                       out.f1 <<- in.f1;
                   }
               }
           }
           when (!counteq1) par {
               out.f1 <<- e;
               count--;
               counteq1 <- (count == 2);
           }
        \end{lstlisting}
        This version easily met the 10ns timing constraint but at the price
        of more obscure coding.






    \section{uncontrolled input sources}
    
        It is often the case that input data is from a source that has no
        flow control, for example a video camera. In that case pipeline
        signalling cannot extend back past the input module. This is a
        problem if any module in the processing pipeline fails to read an
        input queue in the cycle it becomes available. Buffering does not help
        as any pipeline delays accumulate. It is preferable to use a
        processing pipeline with a higher clock frequency than the input
        data. The pipelines are connected by means of an asynchronous queue
        variable which provides a FIFO between modules whose clock
        frequencies may differ.

    \section{bit access}

        There is usually no reason to access individual bits in a word. If a
        word represents flags then it should be an array of {\bf log}. If
        there are packed fields then it should be a {\bf struct}. If data is
        being handled on a word basis but a single bit is to be examined at
        one point in the code one method is to cast the word into an
        array of single bits and then reference the required bit directly,
        e.g.
        \begin{lstlisting}[gobble=8, language=threepl]
           static(a, "uint:8");
           value(b, "[8]bits:1");

           b := cast(a, "[8]bits:1");

           .
           ... <- .. b[3] ..;  // the 4th bit of 'a'
           .
        \end{lstlisting}
        The alternative is to use masking and shifting in order to access bit
        fields. If the mask word and shift count are constant then the two forms
        will generate equivalent code and are efficient. There are functions in
        stdlib.3pl called getbit(), isset() and isclear() which already do this
        and can be used for single bit access. Note that there is no way to
        assign a single bit of a multibit primitive - the whole datum must be
        assigned with the selected bit changed. If single bits are to be
        assigned the type should be declared as an array of type {\bf log} or of
        {\bf bits:1} and later cast to the primitive multibit type required.
    
    \section{pipeline resets}

        Many library modules have optional reset inputs, which are usually
        value mode. These reset inputs are usually intended for recovery from
        error conditions. It must be recognised that the use of these inputs
        will usually destroy the current state of the pipeline. Modules along
        a pipeline may be buffering zero, one or more data and the exact
        position of successive data in the pipeline is usually not known. If
        an asynchronous reset is to be used it is usually necessary to reset
        the entire pipeline to an initial (empty) state and restart it from
        the input end.



    \section{pipeline conflicts}

        There are two pipeline topological features which have the potential
        to cause trouble. The first is alternate queue paths between a source
        and destination module. The second is queue loops.

        \subsection{alternate queue paths}

            When there are two or more queue paths from a source module to a
            destination module there is the potential for the queues to interact. In
            the simple case where the destination module waits on both inputs and
            then reads both, problems arise if the latency through the paths
            differs. Data in the shorter path will reach the destination module
            first and that path will then block. The point of the block will move
            back until all the queue buffers in the path are full. The longer path
            will in the meantime be flowing but not blocking. When the short path
            fills and blocks the source the longer path will cease to read any more
            data from the source and so will start to empty from the input end.
            Eventually data in the longer path will reach the destination module at
            which point the destination will start to read both queues. After a while
            the short path will unblock the source so that data flows again into
            both paths, but soon the data gap in the long path will reach the
            destination which will block again. This is an oversimplified view and
            in practice the queue interactions can lead to a complicated disjointed
            data flow. The data will always be correct and in step and the pipeline
            will not deadlock, but data processing will be sporadic and the
            throughput will be less than it could be.

            The simple solution is to put a buffer (queue) in the shorter
            path. This will take up the excess data and keep the source
            module from blocking, thus maintaining the flow to the longer
            path. The minimum buffer size should be the difference in latency
            between the two paths. A queue variable whose buffer size has been
            increased by setting attribute {\bf depth()} or
            {\bf mindepth()} can be used for this purpose.

            This phenomenon has not been analysed. It can probably be
            modelled using Petri nets.

        \subsection{queue loops}

            A closed queue loop should only be used after careful
            consideration as it can easily lead to deadlock. It is possible
            to set up a loop if either reading of the inputs is conditional
            at the point at which the loop feeds back, or else if modules
            in the loop are initialised so that the pipeline starts already
            in the steady state condition.
